\section{A differentiable EdgeNN example}
\label{app:edge-nn}
This example expresses a shared edge-local neural network followed by a
destination-wise sum. Positions and node features are ordinary
\code{torch.Tensor} objects; the weights are ordinary \code{torch.nn.Parameter}
objects. No native tensor wrapper or handwritten backward is required.
For a radius relation, the program computes
$$
\mathcal{N}(i)=\{j\ne i:\|p_j-p_i\|_2\le r\},\qquad
y_i=\sum_{j\in\mathcal{N}(i)}\left[
W_2\,\operatorname{ReLU}\!\left(W_1[p_j-p_i\,;\,x_j]+b_1\right)+b_2\right].
$$
Here the final bias is part of each edge message, inside the sum.
The feature widths are 8 for node inputs, 3 for displacement, 16 for the
hidden layer and 4 for each output. The reducer produces a zero vector for
a destination with no neighbors.

\subsection{Define the message and differentiate the result}
\code{tg.nn.trace} captures the supported Torch module without replacing its
parameters. Its multiple tensor arguments are concatenated along the feature
axis: \code{edge.displacement} supplies the first three features and
\code{src.x} the remaining eight. \code{reducer = tg.sum()} specifies aggregation
once on the class, rather than inside the edge function.
\lstinputlisting[language=Python,firstline=2,lastline=30]{examples/edge_nn.py}

The call returns one four-component vector per destination. The loss yields
gradients for node features, positions and every MLP weight and bias through
\code{torch.autograd.grad}; \code{loss.backward()} is the usual alternative
when accumulated \code{.grad} fields are wanted. The position derivative acts
through the displacement of selected edges. Discrete radius membership is
not differentiated: moving across the cutoff changes the selected relation.

\subsection{Check against an explicit Torch reference}
The following continuation materializes edges only for validation, applies
the original Torch MLP, and aggregates with \code{index\_add}. Both forward
values and all gradients are compared on the same selected relation.
This isolates message, reducer and differentiation correctness; it is not an
independent validation of the radius-neighbor search.
\lstinputlisting[language=Python,firstline=32]{examples/edge_nn.py}

On the supported CUDA EdgeNN path, Tiga fuses edge evaluation and reduction
into tiled execution and recomputes local activations in the VJP, avoiding
a global edge-by-hidden-feature activation tensor. CPU execution remains
useful for functional checking but does not establish that CUDA performance
property. Unsupported module operations are outside this example's capture
contract. The complete runnable file is
\href{https://github.com/walkerchi/tiga-lang-paper/blob/main/examples/edge_nn.py}{\code{examples/edge\_nn.py}}; execute it
from the paper directory with Tiga and the appropriate Torch installation:
\begin{lstlisting}[language=TigaShell]
python examples/edge_nn.py
\end{lstlisting}
